% Imporditud: tools/import_eksperimendid.py
% Allikas: Eesti füüsikaolümpiaadi eksperimendi ülesannete
%   kogu (Rael Kalda, 2023), ülesanne Ü142, lahendus L142.
\setAuthor{EFO žürii}
\setRound{lõppvoor}
\setYear{2013}
\setNumber{P E2}
\setDifficulty{4}
\setTopic{dünaamika}
\setVahendid{klots, 100-grammise massiga koormis, mõõtejoonlaud, kumminiit}
\setPunktid{12}
\setAllikas{Eksperimendi kogu Ü142, lahendus lk 106}
\setLahendusseis{toores}

\prob{Kumminiit}
Määrake 15 cm võrra väljavenitatud kumminiidi energia.
\textit{Märkus:} 10 cm-st suuremate deformatsioonide korral ei ole kumminiidi pikenemine võrdeline kumminiidile mõjuva jõuga.

\hint

\solu
Teeme kumminiidist ragulka ja paneme selle abil liikuma klotsi. Fikseerime klotsi poolt läbitud teepikkuse, mõõdame hõõrdejõu ja arvutame töö, mis võrdub väljavenitatud kumminiidi energiaga. Taatleme kumminiidi dünamomeetrina. Paneme 100 g vihi kumminiidi otsa ja mõõdame kumminiidi pikenemise 1 N suuruse jõu mõjul. Veame kumminiidi otsas klotsi ühtlase kiirusega mööda lauda ja mõõdame klotsile mõjuva hõõrdejõu. Laseme klotsi liikuma ja mõõdame teepikkuse. (Vähemalt 3 korda) Paneme klotsile koormise ja kordame katset. Arvutame töö ja energia keskmise väärtuse.
\probend
